% !PRINTED 421
\weberpage{421}

Now if $s$ is an arbitrary number coprime to $n$, we can determine $h$ so that
\[
h \equiv s \ (\mathrm{mod}\ n), \quad h \equiv 1 \ (\mathrm{mod}\ a)
\]
and accordingly it follows from (9) that
\[
\varphi(r^s, \alpha) = 0,
\]
that is, equation (7) is satisfied for \gesperrt{all} roots of $X_n$. Hence $\varphi(x,\alpha)$ is divisible by $X_n$, and since on the other hand $X_n$ was assumed to be divisible by $\varphi(x,\alpha)$, it must be identical with $X_n$ up to a numerical factor.

\begin{center}
§. 135.

\gesperrt{The Discriminant of the Cyclotomic Equation.}
\end{center}

The roots of unity of degree $n$ can be represented in transcendental form by
\weq{1}{r = e^{\frac{2\pi i k}{n}} = \cos\frac{2\pi k}{n} + i \sin\frac{2\pi k}{n}.}
This root of unity is primitive or not primitive according as $k$ is coprime to $n$ or not. The simplest among the primitive ones is obtained for $k = 1$, namely
\weq{2}{r_0 = e^{\frac{2\pi i}{n}} = \cos\frac{2\pi}{n} + i \sin\frac{2\pi}{n} = a + bi.}
The angle $2\pi : n$ is the $n$th part of the whole circumference of the circle. The various points of division, which one obtains by laying off this angle on the circumference from an arbitrary initial point, determine the regular $n$-gon inscribed in the circle.

% Registered Bd. I scan, PDF page 447 / printed p. 421, rendered at 600 dpi.
% Exact source crop x=100,y=3150,w=810,h=550; SHA-256 27F4BCA9D93A7A0A2DED6313E3840AD29B999168A364F0FDE053D0091B6B361B.
\weberfigureimage{weber-b1-fig28.png}{0.28}

The points of division can be constructed if one knows $r$, and with it $a$ and $b$ (or even only one of these two quantities). In particular, for the side $s$ of the regular $n$-gon, if the radius $c$ of the circumscribed circle is taken equal to 1, we obtain
\[
s = 2 \sin\frac{\pi}{n} = \sqrt{2\left(1 - \cos\frac{2\pi}{n}\right)}.
\]
The remaining quantities $r$ stand in the same relation to the other points of division. On account of this geometric
%CONT>>
